% ---------------------------------------------------------------------------
% GNU Tochnog -> Tochnog Professional convergence: cycle report and plan.
% Build:  pdflatex convergence-report.tex   (run twice for the TOC)
% ---------------------------------------------------------------------------
\documentclass[11pt,a4paper]{report}

\usepackage[utf8]{inputenc}
\usepackage[T1]{fontenc}
\usepackage{lmodern}
\usepackage[margin=2.5cm]{geometry}
\usepackage{graphicx}
\usepackage{booktabs}
\usepackage{longtable}
\usepackage{array}
\usepackage{xcolor}
\usepackage{hyperref}

\hypersetup{
  colorlinks=true,
  linkcolor=blue!60!black,
  urlcolor=blue!60!black,
  citecolor=blue!60!black,
  pdftitle={GNU Tochnog to Tochnog Professional Convergence},
  pdfsubject={Convergence cycle report and post-convergence improvement plan},
  pdfauthor={[Author name]}
}

% Literal monospace for record/keyword/file names (handles underscores safely).
\newcommand{\code}[1]{\texttt{#1}}
\newcolumntype{L}[1]{>{\raggedright\arraybackslash}p{#1}}
\renewcommand{\arraystretch}{1.12}
\setlength{\parskip}{2pt}
\setlength{\emergencystretch}{2em}

\begin{document}

% ---------------------------------------------------------------------------
\begin{titlepage}
\centering
\vspace*{2.2cm}
{\Huge\bfseries GNU Tochnog $\rightarrow$ Tochnog Professional\\[0.3em]
Convergence\par}
\vspace{0.9cm}
{\Large\bfseries Cycle Report and Post-Convergence\\[0.2em] Improvement Plan\par}
\vspace{1.8cm}
{\large Prepared for the original Tochnog GNU authors\par}
\vspace{0.4cm}
{\large \today\par}
\vfill
{\large [Author name]\par}
\vspace{1.2cm}
{\small
\begin{tabular}{@{}ll@{}}
\textbf{Branch:}   & \code{dev/\allowbreak masin-kernel-surface} at \code{e9c40c9} \\[0.25em]
\textbf{Corpora:}  & sfnet 197, validation-suite 233, internal suite 16 \\[0.25em]
\textbf{Result:}   & \textbf{420/430 combined (97.7\%)}; validation-suite 233/233 \\[0.25em]
\textbf{Sources:}  & \code{ProjectDocs/\allowbreak CONVERGENCE-CYCLE-REPORT.md}, \\
                   & \code{ProjectDocs/\allowbreak IMPROVEMENT-ANALYSIS.md}, \\
                   & \code{ProjectDocs/\allowbreak BACKLOG-OPEN-ISSUES.md}, \\
                   & \code{ProjectDocs/\allowbreak CORPUS-RECONSTRUCTION.md}, \\
                   & \code{ProjectDocs/\allowbreak CORPUS-STALE-TARGETS.md} \\
\end{tabular}
}
\end{titlepage}

% ---------------------------------------------------------------------------
\chapter*{How to read this report}
\addcontentsline{toc}{chapter}{How to read this report}

This report has two audiences: the original Tochnog GNU authors, who know the
code and the physics, and readers who have never seen the project and need the
story before the detail. The technical body is unchanged in substance; the
front matter, the background, and the glossary were added so that a newcomer
can follow it.

\begin{itemize}
  \item \textbf{In a hurry?} Read \emph{Key results in one page} (next page)
        and the \emph{Executive summary}. That is the whole outcome in two
        pages.
  \item \textbf{New to the project?} Read Part~I: the introduction explains
        what Tochnog is, why there are two versions, and what ``convergence''
        means; the glossary defines the domain terms; the context and method
        chapters state the precise scope and rules.
  \item \textbf{Looking for a specific change?} Part~II lists the closed work
        items with their commits and acceptance evidence. Part~III has the
        metrics; Part~IV the open backlog; Part~V the improvement plan.
  \item \textbf{Want the evidence?} The appendices document the corpus
        reconstruction and the reproducibility setup. The per-test ledger is
        \code{ProjectDocs/\allowbreak SEGUIMIENTO-\allowbreak CONVERGENCIA.md}.
\end{itemize}

\textbf{Conventions.} \code{PASS} means the process exited with \code{rc=0}
\emph{and} met the test's own \code{target\_\allowbreak value}; \code{RUNFAIL} means it ran
but missed the target (or timed out); \code{PARSE} means the input was rejected
at parse time. File, record and keyword names are set in \code{monospace}.
Commit hashes identify the exact change on the branch named on the title page.
``The Professional'' always means the Tochnog Professional binary used as the
behavioural reference; ``the fork'' means the GNU line under work.

% ---------------------------------------------------------------------------
\chapter*{Key results in one page}
\addcontentsline{toc}{chapter}{Key results in one page}

\textbf{The question.} The GNU fork of Tochnog and Tochnog Professional share
the same input language and the same physical models, but over the years they
drifted: the same input can give different answers on the two. This cycle asked
\emph{why}, fixed the GNU side where the difference was a genuine defect, and
proved each fix against the Professional binary.

\textbf{The result.} On the two reproducible test corpora:

\begin{center}
\begin{tabular}{@{}L{5.4cm}rrrr@{}}
\toprule
Corpus & Tests & PASS & RUNFAIL & PARSE \\
\midrule
sfnet 197 (in-repo test set)          & 197 & 187 & 9 & 1 \\
validation-suite 233 (reconstructed)  & 233 & \textbf{233} & 0 & 0 \\
\textbf{Combined}                     & \textbf{430} & \textbf{420 (97.7\%)} & \textbf{9} & \textbf{1} \\
Internal build suite                  & 16 & 16 & 0 & 0 \\
\bottomrule
\end{tabular}
\end{center}

\textbf{What was fixed.} A set of genuine defects in the GNU fork, each
validated against the corpora: the Bi-CG direct-LU retry band width (which
silently corrupted coupled consolidation), a missing \code{else} in the
local-refinement dispatch, spurious \code{node\_\allowbreak dof} requests that aborted
mesh-only models, an inverted groundflow velocity sign, a record-length limit
that made a documented camclay form unreachable, and the nodal-stress recovery
operator (the Professional assigns an absolute weighted average; the fork
accumulated increments). The full list is in Part~II.

\textbf{What it means.} A user can now run the same input on the GNU fork and
expect the same behaviour as the Professional on the great majority of models.
No remaining failure is a known source defect: every one is a stale or
malformed test input, an out-of-scope model or infrastructure limit, a narrow
post-processing divergence, or a performance cost (Part~IV).

\textbf{How it was done.} Evidence before change: every fix started from a
reproduced symptom and a reproduced root cause, and was accepted only after a
full corpus A/B showed exactly the intended change with zero regressions.
Hypotheses that the measurements refuted were recorded and \emph{not} acted on
(Part~II, \emph{Refuted hypotheses}).

\textbf{A recovered corpus.} The validation suite's inputs had been deleted;
they were reconstructed from archived \code{.dbs} database dumps, with the
recovery validated against every known ground-truth pair
(Appendix~\ref{app:reconstruction}).

\textbf{Where to go next.} Part~V is the post-convergence improvement plan:
solver robustness, reproducible corpora, and the remaining performance work.
These are proposals, not commitments.

% ---------------------------------------------------------------------------
\chapter*{Executive summary}
\addcontentsline{toc}{chapter}{Executive summary}

This report documents a convergence cycle whose goal was to align the behaviour
of the GNU fork of Tochnog (the sfnet 2014 sources) with Tochnog Professional,
which served as the behavioural reference. The work was driven entirely by
measurement: every candidate fix began from a reproduced symptom and a reproduced
root cause, and was accepted only after a full A/B run of both test corpora
showed exactly the intended classification change with zero regressions.

The cycle closed 24 work items, fixed a set of genuine defects in the GNU fork,
and---equally important---refuted a number of plausible hypotheses before any
speculative edit was made. It also recovered a deleted validation corpus from
archived \code{.dbs} database dumps, using a derivation that was validated
against every known ground-truth pair.

The measured result is \textbf{420 of 430 corpus tests passing (97.7\%)}: sfnet
197 is at \textbf{187 PASS / 9 RUNFAIL / 1 PARSE}, and the reconstructed
validation-suite 233 is at \textbf{233 PASS / 0 RUNFAIL / 0 PARSE}. The internal
build suite remains \textbf{16/16}. No open issue is rated CRITICAL or HIGH:
every remaining failure is a stale or malformed test input, an out-of-scope
model or infrastructure limit, a narrow post-processing divergence, or a
performance cost.

The report is organised in five parts. Part~I states the background (what
Tochnog is and what convergence means), a glossary for newcomers, the context,
and the method.
Part~II presents the 24 closed items with their commits and acceptance evidence,
the mechanisms discovered, and the hypotheses that were refuted. Part~III reports
the results and the current status. Part~IV summarises the open backlog by
severity. Part~V presents the post-convergence improvement plan with its effort
and benefit. The appendices document the corpus-reconstruction method and the
reproducibility notes.

% ---------------------------------------------------------------------------
\tableofcontents

% ===========================================================================
\part{Background and Method}

\chapter{Introduction and background}
\label{ch:intro}

\section{What Tochnog is}

Tochnog is a finite-element program: it solves differential equations over a
mesh by dividing the domain into elements, evaluating the governing physics at
integration points inside each element, assembling a global system, and solving
it for the unknown degrees of freedom. Tochnog's home ground is solid and
geotechnical mechanics (elasticity, plasticity, hypoplasticity, consolidation,
contact and interfaces), with heat conduction, groundwater flow, and coupled
thermo-hydro-mechanical problems alongside. A model is described in a
plain-text input file made of keyword records, so the same input language
drives every problem.

\section{Two lines of the same code}
\label{sec:two-lines}

The program exists in two lines that share a common ancestry and the same input
language:

\begin{itemize}
  \item \textbf{The GNU fork} --- the open-source line, whose last SourceForge
        (sf.net) release is from January 2014. This is the line under work. Its
        test inputs are the ``sfnet'' corpus referred to throughout.
  \item \textbf{Tochnog Professional} --- the proprietary line, maintained on
        the Professional side, which documents a much larger keyword set and
        serves here as the \emph{behavioural reference}.
\end{itemize}

Because the two share the input language and the physics, a user should be able
to take a model written for one, run it on the other, and get the same answer.
Over the years the two lines drifted apart in a number of behaviours, and that
drift is the problem this report is about.

\section{What ``convergence'' means here}
\label{sec:intro-convergence}

Convergence in this report is \emph{behavioural}: close the tests the GNU fork
fails, and for every failure that remains, establish an evidence-backed reason
(source defect, deliberate divergence, stale input or target, or out of scope)
instead of leaving it unexplained. It is \emph{not} a claim that the two
binaries produce bit-identical numbers; the goal is that the same input
produces the same physical behaviour, within the tolerance each test declares.

The Professional binary is the arbiter: where the GNU result differs, the
question is whether the GNU behaviour is a defect relative to the reference.
If it is, the GNU side is fixed and the fix is proved by running both corpora
before and after; if it is not, the difference is documented and classified
rather than patched away.

\section{Why this matters}

For a user, the value is portability and trust: an input written for Tochnog
can be moved between the GNU and Professional lines without silently changing
the answer, and a result obtained with the GNU fork can be checked against the
Professional reference. For the project, the value is a clean, honest baseline:
the remaining failures are known and classified, so future work starts from
evidence rather than from a list of unexplained red tests. The cycle also
recovered a validation suite whose inputs had been deleted, turning an
un-runnable set into a reproducible one.

\section{How the problem was attacked, in plain terms}

The effort was evidence-led. For each failing test, the first step was to
reproduce the symptom and then reproduce the \emph{cause} --- ideally by
comparing the GNU run, the reference GNU 2014 build, and the Professional
binary. A change was made only once the cause was measured. It was then kept
only if a full run of both corpora showed that exactly the intended tests
changed and nothing else regressed. Many plausible explanations were measured
and \emph{refuted} before any speculative edit was made; those refutations are
recorded as results in their own right, because acting on them would have
introduced regressions.

The technical chapters that follow state the same scope and method precisely,
for readers who want the exact definitions and rules.

\chapter{Glossary}
\label{ch:glossary}

This chapter defines the terms a reader coming to the project for the first
time is likely not to know. The definitions are deliberately short; the
technical chapters give the precise usage.

\section*{The project and its versions}
\begin{description}
  \item[Finite-element method (FEM)] A numerical method that solves a
        differential equation by splitting the domain into small pieces
        (elements) and approximating the solution within each.
  \item[Tochnog] A finite-element program for solid and geotechnical
        mechanics, heat and groundwater flow, and coupled problems, driven by a
        keyword-based text input.
  \item[GNU fork] The open-source line of Tochnog; its last sf.net release is
        from January 2014. The line under work.
  \item[Tochnog Professional] The proprietary line, used here as the
        behavioural reference.
  \item[sfnet 197] The in-repo test set taken from the sf.net release: 197 test
        inputs.
  \item[validation-suite 233] A reconstructed suite of 233 test inputs; the
        originals had been deleted and were recovered from \code{.dbs} dumps.
  \item[internal suite 16] The repository's own build smoke tests (16).
  \item[Corpus] A collection of test inputs run together as a measurement
        basis; ``combined 430'' is sfnet 197 plus validation-suite 233.
\end{description}

\section*{Running a test}
\begin{description}
  \item[\code{.dat} (input file)] The plain-text model a user writes, as a
        sequence of keyword records.
  \item[\code{.dbs} (database dump)] The file the binary writes at the end of a
        run. It mixes the parsed input records with run-produced state; this is
        what made the lost inputs recoverable.
  \item[PASS / RUNFAIL / PARSE] The three classifications. \code{PASS}: the
        process exits with \code{rc=0} \emph{and} the test's own
        \code{target\_\allowbreak value} is met. \code{RUNFAIL}: it runs but misses
        the target, or times out. \code{PARSE}: the input is rejected while
        being read.
  \item[\code{target\_\allowbreak value}] The expected value declared inside a test;
        the binary checks it and sets the exit code accordingly.
  \item[\code{rc}] The process exit code; \code{rc=0} means success.
  \item[Work unit] One self-contained change: the code, its verification, and
        its documentation, committed together.
  \item[A/B] Running a corpus before and after a change and diffing the
        per-test classification, to prove that only the intended tests moved.
  \item[Regression] A test that passed before a change and fails after it.
\end{description}

\section*{Numerics and mechanics}
\begin{description}
  \item[Degree of freedom (dof)] An unknown of the discrete problem, attached
        to a node or an element (displacement, temperature, pressure, ...).
  \item[Integration point (IP)] A sampling point inside an element where the
        material response is evaluated; the element integral is built from
        these points.
  \item[Quadrature] The numerical rule (a set of points and weights) used to
        integrate over an element. \textbf{Gauss} and \textbf{Lobatto} are two
        such rules with different point placement and accuracy. The
        Professional uses Gauss for the element integral; the historic GNU
        default was Lobatto.
  \item[Bi-CG (biconjugate gradient)] An iterative solver for the linear
        system. On ill-conditioned systems it may fail to converge; Tochnog
        then retries with a direct solver.
  \item[Direct LU / \code{dgbsv}] A banded direct factorisation used as the
        retry when the iterative solver breaks down.
  \item[Preconditioner] A transformation applied to make an iterative solver
        converge faster (for example, a diagonal/Jacobi scaling).
  \item[mpc / tie] A multi-point constraint: a linear relation between two or
        more degrees of freedom (for example, tying two nodes together at an
        interface).
  \item[Nodal stress recovery] Converting stresses computed at integration
        points into smooth values at nodes, typically by a least-squares fit
        per element followed by an average over the adjacent elements.
  \item[Hypoplasticity / Ma\v{s}in models] A family of soil constitutive models
        (Ma\v{s}in 2005/2007) implemented in the fork and validated in this
        cycle.
  \item[Refuted hypothesis] A candidate cause that measurement disproved. It is
        recorded so that it is not acted on later.
\end{description}

\section*{Severity and effort}
The backlog (Part~IV) rates each open issue by severity
(CRITICAL / HIGH / MEDIUM / LOW) and by effort (cheap / medium / large); the
scales are defined at the head of that part. In short: CRITICAL would corrupt
common results, HIGH a documented feature diverges on a plausible model, MEDIUM
a narrow or legacy divergence, and LOW a single test artifact or a cosmetic
issue.

\chapter{Context}
\label{ch:context}

\section{The GNU fork and the Professional reference}

The project under work is a GNU fork of Tochnog, derived from the sfnet 2014
sources. A separate Tochnog Professional distribution (maintained on the
Professional side) provides a well-defined behavioural reference. The fork and
the Professional share the same input language and the same physical models, but
they have drifted in a number of behaviours over time.

The convergence cycle did not aim at bit-identity. Its target was
\emph{behavioural convergence}: to close the corpus tests that the GNU fork
fails, and, for every remaining failure, to establish an evidence-backed
classification (source defect, deliberate divergence, stale input or target, or
out-of-scope) rather than leaving the failure unattributed.

\section{Corpora}

The measurement basis is three test sets:

\begin{itemize}
  \item \textbf{sfnet 197} --- the in-repo sfnet test set
        (\code{external-downloads/\allowbreak sfnet/\allowbreak extracted/\allowbreak test}).
  \item \textbf{validation-suite 233} --- a reconstructed validation suite
        (\code{validation-suite/\allowbreak test-2014}); the original \code{.dat} inputs had
        been deleted and were recovered from archived \code{.dbs} dumps
        (Appendix~\ref{app:reconstruction}).
  \item \textbf{internal suite 16} --- the repository's own build smoke suite.
\end{itemize}

Both corpus directories are untracked/ignored in the repository. The frequently
cited ``363'' corpus exists only in ephemeral \code{/\allowbreak tmp} and is not reproducible
in-repo; ``\textasciitilde207'' was a PASS count, not a corpus size. The honest,
reproducible denominators are therefore \textbf{197} and \textbf{233}, for a
combined \textbf{430}.

\section{What convergence means here}
\label{sec:convergence-def}

Convergence is measured by the corpus classification plus the internal suite:

\begin{description}
  \item[PASS] the process exits with \code{rc=0} \emph{and} the test's own
        \code{target\_\allowbreak value} is met. The binary enforces this in \code{exit\_\allowbreak tn}
        (\code{miscel.cc}).
  \item[RUNFAIL] the process runs but does not meet the target, or times out.
  \item[PARSE] the input is rejected at parse time (e.g.\ an unknown record).
\end{description}

This is a behavioural metric. It is \emph{not} a claim of bit-identity: for the
reconstructed tests, \code{rc=0} does not imply bit-identity with the
pre-incident archive, because the archive-era binary predates many later physics
commits.

\chapter{Method}
\label{ch:method}

The cycle ran as a sequence of narrow, ordered work units under five invariant
rules.

\begin{enumerate}
  \item \textbf{Evidence before change.} Every candidate fix started from a
        measured symptom and a measured root cause. Source changes were made only
        after the cause was reproduced, ideally against both the Professional
        binary and the reference GNU 2014 build.
  \item \textbf{Falsify before change.} A hypothesis was treated as refuted
        unless the measurement supported it. Several scheduled work units were
        \emph{not} implemented because the evidence refuted them
        (Chapter~\ref{ch:refuted}).
  \item \textbf{Per-family, not per-test.} Fixes were grouped by root cause (a
        parser family, a solver path, a quadrature rule) and then validated by a
        full corpus A/B, so that the per-test classification diff could be shown
        to be exactly the intended set with zero regressions.
  \item \textbf{Clean builds only.} All measurements used clean
        \code{build\_\allowbreak safe.sh} builds (\code{--clean} where the filesystem
        required it), with \code{LC\_\allowbreak ALL=C} and a 45\,s per-test corpus timeout;
        scratch worktrees were kept out of the main tree.
  \item \textbf{Documentation in the same work unit.} Each change updated
        \code{SEGUIMIENTO-\allowbreak CONVERGENCIA.md} and the relevant
        \code{manual-user}/\code{manual-developer} pages, per the project's
        mandatory convention.
\end{enumerate}

% ===========================================================================
\part{The Work Done}

\chapter{Closed items}
\label{ch:closed}

This part is the record of what was actually changed. A reader new to the
project does not need every row: the point of the table is that each change is
traceable to a commit and to the evidence that accepted it. The rows are in
chronological order of closure.

The cycle closed 24 work items. Each row below gives the item, its commit(s) on
\code{dev/\allowbreak masin-kernel-surface} (ancestors of \code{e9c40c9}), and the acceptance
evidence that was recorded when the item was closed.

\begin{longtable}{@{}r L{5.0cm} L{2.6cm} L{5.6cm}@{}}
\toprule
\# & Item & Commit(s) & Acceptance evidence \\
\midrule
\endfirsthead
\toprule
\# & Item & Commit(s) & Acceptance evidence \\
\midrule
\endhead
\bottomrule
\endlastfoot
1 & Masin-basic 2005 rescue: separate basic kernel (hypo7/8/9), Eq.~(15)
    m-tensor, \code{ts3} cross-term coefficient, Eq.~(21) structure evolution
    plus the Masin-2007 Eq.~(8) pyknotropy modification
  & \code{d25906f}, \code{5a98958}, \code{c945833}
  & hypo7 \code{sigyy=-231.789} (Pro $-231.809$), hypo9 $-231.789$ (Pro
    $-231.811$), hypo8 $-230.534$ / $s=1.609$ (Pro $-230.588$; target
    $-230.58 \pm 0.1$); all \code{rc=0}; hypo1--4 \code{.dbs} byte-identical
    (gated). \\
2 & \code{dynamic8}: \code{bounda\_\allowbreak sine} layout (\code{end\_\allowbreak time}), \code{*\_\allowbreak sine}
    consumers, per-operation quadrature, Professional nodal recovery,
    scale-relative Bi-CG breakdown test
  & \code{6888875}, \code{814c9ee}, \code{06e175d}, \code{abf9c29}
  & final \code{disx20 = -9.98525437506e-4} vs Pro $-9.985253775415e-4$ (abs
    $6.0e-11$); prescribed sine matches the Pro series to $3.3e-11$. \\
3 & \code{mpc5}: deleted-element matrix / \code{lhside} zeroing
    (energy-consistent released load) + absolute nodal-stress recovery +
    post-solve stress sync
  & \code{7687e0a}, \code{fa71193}, \code{0b23957}, \code{dae06a3}
  & \code{post\_\allowbreak point\_\allowbreak dof 10 -sigyy = 0.5} at $k=1,2,3,5$ (target $0.5 \pm
    1e-3$); full \code{node\_\allowbreak dof} field $=$ Pro to $\sim 5e-17$. \\
4 & \code{delete1-4} + \code{control\_\allowbreak mesh\_\allowbreak delete\_\allowbreak geometry\_\allowbreak stop} family:
    full-matrix delete-factor scaling, frozen deleted-element IP state,
    stop/\code{\_\allowbreak stop\_\allowbreak geometry} keywords
  & \code{7687e0a}, \code{0735028}, \code{a548415}
  & \code{delete2} \code{rc=0} node 2 \code{sigxx=0.8} (Pro $0.8$);
    \code{delete3} \code{rc=0} \code{sigyy=0.600002285218} (target $0.6 \pm
    1e-4$); family field $=$ Pro. \\
5 & \code{validation\_\allowbreak 2} (Professional's own suite): coupled 1-D consolidation
  & \code{6eb399d} (ancestor)
  & \code{post\_\allowbreak point\_\allowbreak dof 1 -pres = -0.634519741323} vs Pro $-0.6345197413352$
    (10 digits); target $-0.63 \pm 0.05$. \\
6 & \code{validation\_\allowbreak 8}: \code{-veln} normal-velocity \code{bounda\_\allowbreak dof} with
    generated mpc ties
  & \code{f244ac5} (ancestor)
  & 50/50 \code{mpc\_\allowbreak node\_\allowbreak number} + 8/8 \code{mpc\_\allowbreak node\_\allowbreak factor} records
    byte-identical to the Pro; \code{rc=0}. \\
7 & Per-operation quadrature + Professional nodal stress recovery
  & \code{06e175d}
  & \code{dynamic8} $6e-11$; \code{beam2d\_\allowbreak 1} corner \code{sigxx=59.5766} (Pro
    $59.5710$); \code{force5} $-9.8071801515$ (Pro $-9.80718$); \code{post4}
    moment $0.666667$. \\
8 & \code{-total\_\allowbreak linear} incremental formulation (manual Pro 2.2.1): reset
    survives the solve
  & \code{f045d04}
  & reset-only probe \code{sigyy=1} $=$ Pro; \code{total3} (the only sfnet
    \code{-total\_\allowbreak linear} test) \code{rc=0} \code{disx=0.5}. \\
9 & Parse/arity family: variable-length
    \code{control\_\allowbreak timestep\_\allowbreak iterations\_\allowbreak automatic} and
    \code{hypo\_\allowbreak intergranularstrain}
  & \code{8d70054}
  & 7 tests closed (family); full corpus A/B. \\
10 & \code{node\_\allowbreak dof} family: spurious \code{node\_\allowbreak dof} requests in \code{top.cc}
     parse and the \code{so.cc} eigen path
  & \code{ddc9a5b}
  & 6 tests closed (examp7/refine4/examp15, eigen1/eigen3/ho\_heat1); 167
    $\rightarrow$ 173 PASS, 0 regressions. \\
11 & Groundflow velocity sign: restore $v = +k\,\mathrm{grad}(h)$ (manual Pro
     8.2)
  & \code{0838457}
  & \code{condif5} \code{rc=0} \code{node\_\allowbreak dof 2 temp = 1.0} (target $1.0$);
    full node vector $=$ Pro (\code{gvelx=+10}). \\
12 & \code{refine\_\allowbreak l} BAR2 fall-through: missing \code{else if} in the
     element-shape dispatch
  & \code{0d0ad0c}
  & \code{refine5} \code{rc=0} \code{post\_\allowbreak point\_\allowbreak dof 0 -temp = 1} (target $1 \pm
    5e-2$); the 2014 reference aborts at the same assert. \\
13 & Bi-CG direct-LU retry band width: \code{diff*nprinc} $\rightarrow$
     \code{(1+diff)*nprinc}
  & \code{6eb399d}
  & \code{examp2} $-0.634519741323$ $=$ Pro (was $-0.7206$); \code{validation\_\allowbreak 2}
    closed. \\
14 & Corpus reconstruction from archived \code{.dbs}: 192 \code{.dat} recovered
  & \code{e983be8}, \code{5c80b60}
  & 148 recovered first pass, plus doc-grounded recoveries; derivation validated
    11/11 known pairs, 72/72 IP counts. \\
15 & Stale sfnet input/target alignment: 8 \code{inertia\_\allowbreak apply -yes},
     \code{matrix8} sign, \code{incnav5} target
  & \code{b83d907} (docs)
  & 174 $\rightarrow$ 184 PASS, 10 RUNFAIL $\rightarrow$ PASS, 0 regressions. \\
16 & Interface stiffness guard + reset scoping/timing: guard the optional
     \code{group\_\allowbreak interface\_\allowbreak materi\_\allowbreak elasti\_\allowbreak stiffness}; reset at its own index;
     reset after the state update
  & \code{f142bdf}, \code{a548415}
  & \code{iface\_\allowbreak condif}, \code{groundflow\_\allowbreak interface}, \code{creset\_\allowbreak geom} closed;
    \code{delete3} reconciled. \\
17 & \code{group\_\allowbreak materi\_\allowbreak plasti\_\allowbreak camclay} variable length 3..4: the legacy 4-value
     form was unreachable
  & \code{91676ab}
  & \code{mc\_\allowbreak pressure\_\allowbreak min}/\code{\_\allowbreak off} \code{rc=0} exact doc targets; 4-value
    form now reachable; corpus classification byte-identical. \\
18 & \code{delete3} regression reconciliation: refine the
     \code{control\_\allowbreak reset\_\allowbreak dof} per-index gate (one-shot vs step-scoped)
  & \code{a548415}
  & \code{delete3} \code{rc=0} \code{0.600002285218}; \code{creset\_\allowbreak geom}
    preserved; sfnet A/B byte-identical. \\
19 & Absolute nodal-stress recovery: the Pro assigns the weighted average, GNU
     accumulated increments
  & \code{0b23957}
  & \code{delete2} $0.8$; delete family + mpc1--6 field $=$ Pro; all 191 sfnet
    \code{.dbs} dumps byte-identical. \\
20 & Stale validation-suite target recalibration + \code{gforce7q4} rebuild
  & \code{c592843}
  & 14 targets recalibrated + \code{gforce7q4} rebuilt as the faithful 2-quad4
    Pro model; 201 $\rightarrow$ 216 PASS, 15 RUNFAIL $\rightarrow$ PASS, 0
    regressions. \\
21 & \code{inertia\_\allowbreak apply -yes} reconstruction of the 15 solver-divergence tests
   & \code{c2f7f6c}
   & 13/15 converge to their stored targets; validation-suite 216 $\rightarrow$
     229 PASS; the 2 residuals carry separate causes (Chapter~\ref{ch:backlog}). \\
22 & Reconstruction residuals: \code{iface\_\allowbreak tangref} target recalibrated to
     the honest per-IP value (4.9998), and \code{iface\_\allowbreak gen\_\allowbreak method\_\allowbreak gen} restored
     to the archive's input/generated element split
   & \code{fbcf144}
   & 2 RUNFAIL $\rightarrow$ PASS; validation-suite 229 $\rightarrow$ 231 PASS, 0
     regressions; sfnet byte-identical. \\
23 & Post-solve stress pass hoisted out of the equilibrium loop (performance;
     improvement~(a))
   & \code{3957809}
   & 1.12--1.25$\times$ on \code{large3}/\code{dynamic8}; mpc5/delete2/delete3
     targets unchanged; corpus classification diff 0; internal 16/16. \\
24 & \code{-follow\_\allowbreak material} reference geometry + initial
     \code{element\_\allowbreak truss\_\allowbreak force} self-stress (\code{elmethod}, \code{numit})
   & \code{c8ee32e}
   & validation-suite 231 $\rightarrow$ 233 real tests \textbf{233/233};
     reconstructed inputs byte-identical to the archived \code{.dbs}; sfnet
     unchanged. \\
\end{longtable}

\chapter{Key mechanisms discovered}
\label{ch:mechanisms}

\section{The .dbs reconstruction method}
\label{sec:dbs-method}

The validation-suite \code{.dat} inputs were lost to a too-wide \code{rm -f
*.dat} (2026-08-29) and the directory is not under version control. They were
recovered from the archived \code{.dbs} database dumps, which the binary writes
at the end of a run and which mix the parsed input records (canonicalised,
comment-free, reordered) with run-produced records.

The method keeps the input records and drops the derived ones. The core of the
drop-list is the 45 record names that appear in the 41 known \code{.dat}/\code{.dbs}
pairs in \emph{no} \code{.dat} (state/results such as \code{node\_\allowbreak dof*},
\code{element\_\allowbreak dof*}, \code{node\_\allowbreak rhside*}, \code{element\_\allowbreak matrix*}, \code{dof\_\allowbreak *},
\code{time\_\allowbreak *}, \code{global\_\allowbreak solver\_\allowbreak *}, \code{post\_\allowbreak *}), extended with the
run-produced records seen in the wider dump set; 103 names are dropped in total.
The input \code{node\_\allowbreak dof} is recovered from \code{node\_\allowbreak dof\_\allowbreak start\_\allowbreak refined}. The
one lost scalar is derived rather than invented (Appendix~\ref{app:reconstruction}).

A critical corollary: \code{inertia\_\allowbreak apply} is never echoed in a \code{.dbs} (0 of
430 dumps). Its absence from an archived dump is therefore \emph{not} evidence
that the input lacked it. This directly enabled the recovery of the 15
solver-divergence tests.

\section{Per-operation quadrature}
\label{sec:quadrature}

The Professional does not use a single global quadrature. The quadrature is a
property of the element integral, and the nodal stress recovery is a separate
operator:

\begin{center}
\begin{tabular}{@{}L{6.2cm}L{4.4cm}L{4.0cm}@{}}
\toprule
Operation & Professional & GNU (aligned) \\
\midrule
Element stiffness, internal force, body load (quad4/hex8) & $2\times2$ / $2\times2\times2$ Gauss & Gauss \\
Element stiffness (quad9/hex27) & $3\times3$ / $3\times3\times3$ Gauss & Gauss \\
Mass matrix & row-sum lumped (quadrature-invariant) & unchanged \\
Edge / surface load, interfaces & boundary integration; interfaces keep nodal (Lobatto) & unchanged \\
Stress recovery (IP $\rightarrow$ node) & per-element least-squares fit + arithmetic mean & Professional operator \\
\bottomrule
\end{tabular}
\end{center}

The historical GNU default was Lobatto for the element integral plus an
\code{h}-weighted lumped recovery, which conflated the two. Manual 6.623's
``inertia $\Rightarrow$ lobatto'' clause was measured to describe the
\emph{mass} (which is row-sum lumped and quadrature-invariant), not the
stiffness; the 02-08-2026 binary contradicts it for the stiffness. The alignment
is in \code{06e175d}.

\section{The Professional's nodal stress recovery}
\label{sec:stress-recovery}

The \code{.dbs} dumps expose the operator chain
\code{element\_\allowbreak intpnt\_\allowbreak dof} $\rightarrow$ \code{element\_\allowbreak dof} $\rightarrow$
\code{node\_\allowbreak dof}, with \code{node\_\allowbreak lhside} as the assembled weight. The measured
operators are:

\begin{itemize}
  \item quad4: bilinear interpolation $1, x, y, xy$ (residual vs
        \code{element\_\allowbreak dof} $1.4e-7$);
  \item quad9: uniform-weight linear least squares $1, x, y$ ($1.2e-12$);
  \item hex8: uniform-weight linear least squares $1, x, y, z$ ($1.0e-8$);
  \item hex27: assumed linear least squares (not independently verified).
\end{itemize}

The operator is fixed in the element natural coordinates and uses \emph{uniform}
weights (the Gauss weights give a different operator). Assembly is the
\emph{arithmetic mean} of the adjacent elements' \code{element\_\allowbreak dof}.

The 2026-09-13 alignment fixed the weights, but the GNU still \emph{accumulated}
increments (\code{node\_\allowbreak dof\_\allowbreak new += node\_\allowbreak rhside/\allowbreak node\_\allowbreak lhside}). At a node shared by
two elements whose increments cancel, the reported stress stayed at its
step-start value (\code{delete2} reported $0.6$ where the Pro reports $0.8$). The
Professional \emph{assigns} the absolute weighted average; the fix
(\code{0b23957}) makes the stress dofs assign the ratio, with a gate that keeps
the historic incremental update for convected stress dofs and for models with an
independent strain unknown.

\section{The mpc5 / delete energy consistency}
\label{sec:mpc5}

Two distinct mechanisms closed the delete/mpc family.

\textbf{Delete-factor scaling} (\code{7687e0a} $\rightarrow$ \code{fa71193}). A
partially deleted element stores \code{element\_\allowbreak delete\_\allowbreak factor} plus a frozen
\code{element\_\allowbreak rhside\_\allowbreak delete} (the released-stress internal force). The
Professional stores \emph{no} \code{element\_\allowbreak matrix\_\allowbreak delete} record. The first fix
scaled the full matrix by the factor; the final fix recognises that the frozen
released load is a \emph{constant}, so its tangent is zero: with
\code{ELEMENT\_\allowbreak DELETE\_\allowbreak FACTOR} active the element matrix and \code{element\_\allowbreak lhside}
are set to \emph{zero} (not scaled). This restores the one-pass Newton
cancellation $A = dt\,K = dt\,B^{\mathsf{T}} C B$ of the staggered scheme. The
historic $\mathit{factor}\cdot K$ left a residual stiffness that responded to
strain while the internal force was frozen, so the released load relaxed
geometrically (\code{mpc5} contracted at rate exactly $0.5$, reaching the target
only at $\sim$12 iterations instead of 2).

\textbf{Post-solve stress sync} (\code{dae06a3}). The \code{materi\_\allowbreak stress} dofs
are non-principal: they are updated \emph{after} the solve from the \emph{pre-solve}
constitutive stress, so the reported nodal stress trailed the solved velocity by
one iteration (\code{mpc5} reported $1.0$ at $k=1$; the Pro reports $0.5$). The
fix adds a stress-only \code{element\_\allowbreak loop()} pass after \code{solve()} and before
\code{parallel\_\allowbreak new\_\allowbreak dof\_\allowbreak diagonal()}, re-evaluating the constitutive stress from
the solved velocity. This is \emph{not} algebraically the same as iteration
$k+1$: it adds a constitutive evaluation, not a second momentum solve.

\section{The inertia-apply default gotcha}
\label{sec:inertia}

Commit \code{0560dcf} aligned the GNU \code{options\_\allowbreak inertia} default with the
Professional: an absent \code{inertia\_\allowbreak apply} now means \code{-no} (manual Pro
6.785), where the historic GNU default was \code{-yes}. This deliberate alignment
fixed 13 static corpus tests but silently changed the behaviour of legacy models
that relied on the old default:

\begin{itemize}
  \item 8 sfnet inputs plus \code{incnav3} and \code{examp2} needed an explicit
        \code{inertia\_\allowbreak apply -yes} (the Pro's own copies carry it);
  \item 15 validation-suite models needed it too. Those models have a material
        density but insufficient Dirichlet constraints, so the legacy \code{-yes}
        supplied a \code{rho/\allowbreak dtime} lumped mass diagonal that regularised an
        otherwise \emph{singular} static system. Without it the matrix has
        genuine rigid-body modes, the Bi-CG breaks down and the direct-LU retry
        diverges.
\end{itemize}

The onset was isolated by a clean-build bisect: \code{0560dcf} (09-05 07:06)
fails all 15, while its predecessor and the quadrature commit \code{06e175d} do
not. The \code{-no} default is a deliberate Professional alignment and was
\emph{not} reverted; the fix is the explicit record in the reconstructed inputs.

\section{The honest solver as an explanation for ``divergences''}
\label{sec:honest-solver}

\code{so\_\allowbreak bicg.cc} previously reported success on three non-convergence exits. The
current honest criteria report the true residual, so several tests that ``passed''
on the archive-era binary now correctly report divergence. This is a deliberate
improvement, and it is why some archived targets were met only by a dishonest
success. It also exposed a latent defect: the direct-LU retry band width used
\code{diff*nprinc} instead of \code{(1+diff)*nprinc}, so the first-to-last dof
coupling of adjacent nodes wrote outside \code{mat}/\code{ldmat} and silently
corrupted the retry (\code{6eb399d}).

\section{The \code{-follow\_\allowbreak material} reference geometry}
\label{sec:follow-material}

Under \code{-follow\_\allowbreak material} the mesh follows the material, so the
\code{node} record in a \code{.dbs} dump is the \emph{deformed} geometry, while
the input reference is preserved in \code{node\_\allowbreak start\_\allowbreak refined}
(\code{top.cc} copies the input \code{NODE} there at the start of the run). The
reconstruction had used the deformed \code{node} as input, which is why
\code{elmethod} and \code{numit} would not reproduce their archived targets. The
second lost record was the initial \code{element\_\allowbreak truss\_\allowbreak force}
self-stress, which the dump stores only in its relaxed final state
($\approx 0$). With the reference geometry restored and the initial force
derived from the documented relaxation law
($F_{\mathrm{init}} = k\,(L_{\mathrm{init}} - L_{\mathrm{final}})$, with
$k = EA/L_{\mathrm{init}}$), the reconstructed inputs reproduce the archived
\code{.dbs} byte-for-byte. This is a corpus-reconstruction mechanism, not a
source change.

\chapter{Refuted hypotheses}
\label{ch:refuted}

The cycle's discipline was to treat every scheduled fix as a hypothesis and to
record the refutations as first-class results. The following hypotheses were
refuted or corrected, and several of them were load-bearing: acting on any of
them without the measurement would have introduced a regression.

\begin{itemize}
  \item \textbf{Masin-basic R2 (``reference geometry refuted'')} was a \emph{false
        negative}: the parked kernel had a transcription defect (Eq.~15 used
        $T_{\mathrm{hat}}{:}T_{\mathrm{hat}}^{*}$ instead of
        $T_{\mathrm{hat}}{:}T_{\mathrm{hat}}$), so the earlier ``R2 does not
        improve the residual'' verdict was measured with the wrong kernel. With
        the kernel fixed, R2 is part of the hypo7/8/9 closure.
  \item \textbf{``\code{updated\_\allowbreak linear} log vs linear''} as the cause of the Masin
        stress residual was refuted by the R2 measurement: the strain measure
        matches the Pro (\code{eptyy = -0.300000000}) but the stress residual
        does not improve systematically.
  \item \textbf{Monolithic saddle-point mixed solve} (\code{mpc5} candidate (a))
        was refuted as structurally inconsistent for shared nodes in the current
        tie scheme (mpc1's shared $\sigma_2$ cancels to $1.667$, not $2.0$).
        Candidate (b)---post-solve re-evaluation---was implemented instead.
  \item \textbf{``Candidate (b) is algebraically identical to iteration $k+1$''}
        was corrected: the \emph{value} sequence coincides, but the
        \emph{mechanism} (a constitutive evaluation, not a solve) is distinct and
        does sync the $k=1$ report.
  \item \textbf{``\code{examp6} is pre-existing''} was corrected to a regression
        from the \code{ELEMENT\_\allowbreak DOF} stress-slot fix \code{8f16dc0}, compounded by
        \code{0560dcf}; still out-of-scope because the target is tied to a
        residue-driven adaptive-remesh node numbering.
  \item \textbf{``\code{examp24} cost is mesh generation''} was corrected to
        $O(N^2)$ mesh \emph{mapping} (\code{map.cc} \code{point\_\allowbreak el}/\code{matrix\_\allowbreak ab},
        29M \code{point\_\allowbreak el} calls per cycle).
  \item \textbf{\code{validation\_\allowbreak 2}'s $-0.688985$} was refuted as a corrupt-retry
        artifact; the honest value is $-0.634519741323 = $ Pro.
  \item \textbf{\code{validation\_\allowbreak 8}'s ``not deducible without more data''} was
        refuted: \code{-veln} was implemented and reproduces the Pro's 50/50 +
        8/8 MPC records.
  \item \textbf{The 15 validation-suite solver divergences} were first attributed
        to the honest Bi-CG criteria and the quadrature commit; both were refuted
        by the bisect, and the true cause (\code{inertia\_\allowbreak apply}) was found.
  \item \textbf{The corpus basis itself} was corrected: ``363'' is ephemeral
        \code{/\allowbreak tmp}, ``\textasciitilde207'' was a PASS count; the reproducible
        denominators are 197 and 233.
\end{itemize}

The pattern is consistent: the fixes that survived were the ones whose root cause
had been measured against a reference binary, and the refutations prevented
speculative edits that would have regressed the corpora.

% ===========================================================================
\part{Results and Current Status}

\chapter{Metrics}
\label{ch:metrics}

\section{Start and end of the cycle}

\begin{center}
\begin{tabular}{@{}L{3.6cm}L{4.4cm}L{4.4cm}r@{}}
\toprule
Corpus & Start of cycle (PASS) & End of cycle (PASS) & Net \\
\midrule
sfnet 197 & 160 / 197 & \textbf{187 / 197} & $+27$ \\
validation-suite 233 & 201 / 233\textsuperscript{\dag} & \textbf{233 / 233} & $+32$ \\
Combined 430 & 361 / 430 (84.0\%)\textsuperscript{\ddag} & \textbf{420 / 430 (97.7\%)} & $+59$ \\
Internal suite & 16 / 16 & \textbf{16 / 16} & --- \\
\bottomrule
\end{tabular}
\end{center}

\textsuperscript{\dag}~The validation-suite start figure is the first full
classification after the corpus was reconstructed from the archived \code{.dbs}
dumps (the \code{.dat} inputs had been deleted; Appendix~\ref{app:reconstruction}).
Before reconstruction the corpus was effectively 0 runnable tests.

\textsuperscript{\ddag}~Sum of the per-corpus start figures (160 and 201). The
source document's combined-start cell reads 388 (90.2\%), which equals the sfnet
\emph{end} (187) plus the validation start (201); the per-corpus starts sum to
361. The combined \emph{end} of 420/430 (97.7\%) is consistent with both corpora's
ends (187 + 233).

\section{Convergence table}

\begin{center}
\begin{tabular}{@{}L{4.4cm}rrrrr@{}}
\toprule
Corpus & Size & PASS & RUNFAIL & PARSE & PASS rate \\
\midrule
sfnet 197 & 197 & 187 & 9 & 1 & 94.9\% \\
validation-suite 233 & 233 & \textbf{233} & 0 & 0 & 100\% \\
\textbf{Combined} & \textbf{430} & \textbf{420} & \textbf{9} & \textbf{1} & \textbf{97.7\%} \\
Internal suite & 16 & 16 & 0 & 0 & 100\% \\
\bottomrule
\end{tabular}
\end{center}

\section{Classification caveats}

\begin{itemize}
  \item \code{PASS} is \code{rc=0} with the test's own \code{target\_\allowbreak value} met,
        which the binary enforces in \code{exit\_\allowbreak tn} (\code{miscel.cc}).
  \item For the reconstructed tests, \code{rc=0} does not imply bit-identity with
        the pre-incident archive: the archive-era binary predates many later
        physics commits, so a faithful reconstruction can still drift
        numerically.
  \item The final sfnet bucket is \code{9 RUNFAIL /\allowbreak  1 PARSE} rather than
        \code{10 RUNFAIL /\allowbreak  0} only because \code{pruebavel} is classified as
        \code{PARSE} (unknown-record error); the PASS count is unchanged.
  \item The validation-suite directory holds one generated dump file
        (\code{print\_\allowbreak mesh\_\allowbreak dof.dat}) that is not a test input. The 233
        real tests are all \code{PASS}; the dump is excluded from the
        denominator.
\end{itemize}

% ===========================================================================
\part{Open Issues}

\chapter{Backlog}
\label{ch:backlog}

Every remaining non-PASS corpus test and every residual found during the cycle is
enumerated in \code{ProjectDocs/\allowbreak BACKLOG-OPEN-ISSUES.md}, with symptom, root
cause, evidence, severity, proposed fix and estimated effort.

\section{Severity and effort scales}

\begin{center}
\begin{tabular}{@{}L{2.2cm}L{11.4cm}@{}}
\toprule
Level & Meaning \\
\midrule
\textbf{CRITICAL} & silently corrupts results on common models, crashes on valid input, or blocks a whole feature family with no workaround \\
\textbf{HIGH} & a documented feature diverges from the Professional on a reproducible model a plausible user would hit \\
\textbf{MEDIUM} & a narrow or legacy divergence, a coverage gap, or a performance cost that affects a common configuration \\
\textbf{LOW} & a single test artifact, a stale/malformed input, or a cosmetic/documentation issue \\
\bottomrule
\end{tabular}
\end{center}

\begin{center}
\begin{tabular}{@{}L{2.2cm}L{11.4cm}@{}}
\toprule
Level & Meaning \\
\midrule
\textbf{cheap} & under a day, isolated change \\
\textbf{medium} & days, one subsystem, needs an A/B corpus run \\
\textbf{large} & weeks, cross-cutting refactor or new infrastructure \\
\bottomrule
\end{tabular}
\end{center}

\section{Open issues by severity}

The table below is the current open backlog. Four validation-suite residuals
and the performance cost that were still open at the cycle snapshot have since
been closed (see \emph{Resolved since the cycle snapshot} below); the
validation-suite corpus now has no open failure. The remaining entries are all
sfnet, plus the unexercised alias, the interface sign residual, and the
interface-generation coverage gap.

\begin{longtable}{@{}r L{2.5cm} L{7.4cm} L{1.7cm} L{2.2cm}@{}}
\toprule
Id & Corpus & Item & Severity & Effort \\
\midrule
\endfirsthead
\toprule
Id & Corpus & Item & Severity & Effort \\
\midrule
\endhead
\bottomrule
\endlastfoot
B1 & sfnet & \code{examp6} remesh node numbering & MEDIUM & cheap/large \\
B2 & sfnet & \code{examp14} stale target + runtime & LOW & cheap \\
B3 & sfnet & \code{examp17} stale target & LOW & medium \\
B4 & sfnet & \code{examp24} $O(N^2)$ mesh mapping & LOW & large \\
B5 & sfnet & \code{ho\_\allowbreak mech2} stale p-refinement input & LOW & medium \\
B6 & sfnet & \code{ho\_\allowbreak mech3} model divergence & LOW & large \\
B7 & sfnet & \code{ho\_\allowbreak mech4} missing strain record & LOW & medium \\
B8 & sfnet & \code{petsc1} PETSc infrastructure & LOW & large \\
B9 & sfnet & \code{taylor3} stale input + solver & LOW & medium \\
C1 & sfnet & \code{pruebavel} dead/malformed & LOW & cheap \\
C2 & validation & \code{groundflow\_\allowbreak interface} \code{gvelx} sign & LOW & medium \\
C3 & sfnet & \code{gforce10} 3D section sign & MEDIUM & medium \\
C5 & --- & \code{contact\_\allowbreak target\_\allowbreak geometry} index-0 alias & MEDIUM & medium \\
C6 & validation & \code{iface\_\allowbreak gen*} generation-coverage gap & MEDIUM & cheap \\
\end{longtable}

\section{Resolved since the cycle snapshot}

These items were open when the first version of this report was written and
have since been closed; they are listed here so the history is complete.

\begin{itemize}
  \item \textbf{A1 \code{iface\_\allowbreak tangref}} --- the target was recalibrated to the
        honest per-IP value (4.9998); the archived 10.3533 came from the
        pre-\code{bf74f28} scalar interface assembly. Closed by \code{fbcf144}.
  \item \textbf{A2 \code{iface\_\allowbreak gen\_\allowbreak method\_\allowbreak gen}} --- the reconstruction was
        corrected to the archive's input/generated element split, so the target
        now reads the \emph{generated} \code{element\_\allowbreak geometry} record rather than
        an input tag. Closed by \code{fbcf144}.
  \item \textbf{A3 \code{elmethod} / A4 \code{numit}} --- the \code{-follow\_\allowbreak material}
        reference geometry and the initial \code{element\_\allowbreak truss\_\allowbreak force}
        self-stress were recovered; the reconstructed inputs now reproduce the
        archived \code{.dbs} byte-for-byte. Closed by \code{c8ee32e}
        (Section~\ref{sec:follow-material}).
  \item \textbf{C4 post-solve stress pass cost} --- the stress pass was hoisted
        out of the equilibrium loop, removing the duplicated element pass
        (1.12--1.25$\times$ measured). Closed by \code{3957809}
        (improvement~(a)).
\end{itemize}

\section{Summary}

No open issue is CRITICAL or HIGH, and the validation-suite corpus has no open
failure: its 233 tests all pass. Fourteen items remain open: four MEDIUM and ten
LOW. Every remaining failure is a stale or malformed test, an out-of-scope model
or infrastructure limit, a narrow post-processing divergence, a coverage gap, or
a performance cost. The source defects found during the cycle were fixed and
validated against the corpora.

The remaining sfnet failures are the stale or malformed inputs, the $O(N^2)$
mesh-mapping runtime, the PETSc build infrastructure, and two post-processing
divergences (\code{gforce10}'s 3D section sign, and the interface \code{gvelx}
convention residual). The cross-corpus items left are the unexercised
\code{contact\_\allowbreak target\_\allowbreak geometry} index-0 alias and the
\code{iface\_\allowbreak gen*} generation-coverage gap (the reconstructed inputs
pass via \code{node\_\allowbreak dof} but do not exercise interface generation).

% ===========================================================================
\part{Future Plans}

\chapter{Improvement plan}
\label{ch:improvements}

The improvements below were identified during the cycle. Each states the problem,
the proposed change, the expected benefit, the risk and the estimated effort.
These are proposals, not commitments; none is required for the cycle's results.
The full analysis is in \code{ProjectDocs/\allowbreak IMPROVEMENT-ANALYSIS.md}.

\section{Summary}

\begin{longtable}{@{}r L{5.2cm} L{3.4cm} L{2.2cm} L{2.2cm}@{}}
\toprule
Id & Improvement & Benefit & Risk & Effort \\
\midrule
\endfirsthead
\toprule
Id & Improvement & Benefit & Risk & Effort \\
\midrule
\endhead
\bottomrule
\endlastfoot
a  & Hoist the post-solve stress pass & \textbf{DONE} (1.12--1.25$\times$) & low--moderate & cheap--medium \\
b  & Preconditioner; revisit the mixed solve & Robustness; structural fix & medium--high & medium/large \\
c  & Clean the \code{build\_\allowbreak safe.sh} full-mode file checks & A clean CI gate & moderate & medium \\
d  & Reproducible corpora + stable runner + recalibration tooling & Reproducibility, CI & low--moderate & medium \\
e  & Promote the \code{.dbs} reconstruction tooling & Re-runnable, self-checking recovery & low & medium \\
f1 & Single quadrature-rule helper & One place to change & low & cheap \\
f2 & Principal stress dofs (mixed solve) & Structural correctness & high & large \\
f3 & Reset scoping & One clear semantic & low & cheap \\
f4 & Retire \code{ELEMENT\_\allowbreak MATRIX\_\allowbreak DELETE} & Less dead code & low & cheap \\
f5 & Fix the \code{-BAR3} refine branch & Correct, tested path & low & cheap \\
f6 & Exact band width / bounds assertion & Tighter band; loud failure & low & cheap/medium \\
f7 & Document \code{control\_\allowbreak eigen} as GNU-only & Avoid a phantom work unit & none & cheap \\
\end{longtable}

\section{Item narratives}

\begin{description}
  \item[(a) Post-solve stress pass cost --- \textbf{DONE}.] The post-solve
        nodal-stress sync (\code{dae06a3}) ran a stress-only
        \code{element\_\allowbreak loop()} after \emph{every} equilibrium iteration for
        models with \code{materi\_\allowbreak stress} and the absolute recovery, roughly
        doubling the element-assembly cost for those models. Only the final
        state of a step is reported, so the pass was hoisted out of the
        equilibrium loop (\code{3957809}): it now runs once per step, after the
        loop, and assigns the stress dofs directly via a stress-only variant of
        \code{parallel\_\allowbreak new\_\allowbreak dof\_\allowbreak diagonal}. Measured speedups:
        \code{large3} (784 elements, 220 iterations) 1.23--1.25$\times$,
        \code{dynamic8} (16 elements, 4000 iterations) 1.12--1.21$\times$; the
        small delete/mpc tests are startup-dominated. Targets unchanged
        (mpc5 \code{sigyy=0.5}, delete2 \code{sigxx=0.8}, delete3
        \code{sigyy=0.6}); corpus classification diff 0; internal suite 16/16.

  \item[(b) Solver robustness.] The solver is honest but not robust on the
        genuinely hard systems. Two effects: the honest Bi-CG criteria (correct,
        but ill-conditioned systems fall to the direct-LU retry), and the
        genuinely hard systems (consolidation coupling, the singular static
        systems that relied on the legacy inertia default, and the eulerian
        convection/plasticity family). Two directions, in order of cost: a simple
        diagonal (Jacobi) or block-Jacobi preconditioner to reduce the iteration
        count without changing the equations; and, as a long-term principled fix,
        MINRES or a mixed solve for the staggered family. The latter was refuted
        as structurally inconsistent for shared nodes in the current tie scheme,
        and remains valid only \emph{if} the tie scheme is made
        energy-consistent for shared nodes first (the current
        energy-consistent tie work is the precondition for revisiting it).

  \item[(c) \code{build\_\allowbreak safe.sh} full-mode file-check layer.] After the corpus
        reconstruction the canary \code{msf\_\allowbreak beam2d.dat} exists, so the script
        runs in \textbf{full mode} and executes the Sprint~11 file-check layer.
        That layer asserts expected outputs of files generated during the
        pre-incident sprints; many are stale relative to the current binary, so
        the script exits non-zero even though the FALLO list is identical to the
        baseline. The proposal is to update the stale expected outputs to the
        current honest values (with the same A/B discipline used for the corpus
        target recalibrations), or to move the file-check expectations into a
        maintained golden-file set generated from a known-good binary and
        reviewed like any other target. Keep the \code{rc}-based internal suite
        (16/16) as the primary gate.

  \item[(d) Corpus tooling.] The measurement basis is not reproducible from the
        repository: the ``363'' corpus lives only in ephemeral \code{/\allowbreak tmp}; the
        runner is a \code{/\allowbreak tmp} shell script with ad-hoc classification output;
        and target recalibration is manual. Three pieces are proposed: a
        \code{corpus/\allowbreak } manifest with checksums plus a fetch/extract script (the
        third-party test data has licensing implications, so a fetch script with
        pinned checksums is safer than committing the data); a stable runner at
        \code{scripts/\allowbreak run-corpus.sh} with machine-readable output
        (\code{name}, tab, \code{PASS|RUNFAIL|PARSE}, tab, \code{value}); and
        target-recalibration tooling that proposes the edit while leaving the
        human decision explicit.

  \item[(e) Reusability of the reconstruction method.] The \code{.dbs}
        reconstruction method is well documented but implemented in ad-hoc
        \code{/\allowbreak tmp} scripts. It should be promoted to
        \code{scripts/\allowbreak reconstruct\_\allowbreak dbs.py}: the drop-list as a versioned data file
        (name $\rightarrow$ reason); the \code{number\_\allowbreak of\_\allowbreak integration\_\allowbreak points}
        derivation with the exact-division check; a \textbf{validation mode} that
        runs the derivation against the known \code{.dat}/\code{.dbs} pairs (the
        11/11 exact check) so a regression in the method is caught immediately;
        and a RECOVERED / PARTIAL / UNRECOVERABLE report matching the documented
        classification.

  \item[(f1) Three replicated quadrature defaults.] The element-integral
        quadrature default is replicated in three places: the element default,
        \code{msf\_\allowbreak element\_\allowbreak rule} (\code{calcul\_\allowbreak force.cc}) and \code{mc\_\allowbreak element\_\allowbreak rule}
        (\code{delete.cc}). They were switched to Gauss together in
        \code{06e175d}, but a future change must keep them in sync manually.
        Resolve the rule from a single helper (e.g.\
        \code{element\_\allowbreak quadrature\_\allowbreak rule()}), with the two call sites reading it.

  \item[(f2) Non-principal stress dofs.] The \code{materi\_\allowbreak stress} dofs are
        non-principal (\code{input.cc} sets no \code{dof\_\allowbreak principal}), so they are
        updated after the solve from the pre-solve constitutive stress. The
        one-iteration lag was closed by a post-solve pass (\code{dae06a3}), which
        is a workaround, not the structural fix, and which carries the
        performance cost of item~(a). Revisit the monolithic saddle-point mixed
        solve \emph{after} making the tie scheme energy-consistent for shared
        nodes. If that is not pursued, keep the post-solve pass as the documented
        mechanism and apply item~(a).

  \item[(f3) Reset scoping legacy branch.] \code{control\_\allowbreak reset\_\allowbreak dof} keeps a
        legacy every-step behaviour for resets whose index also carries a
        \code{control\_\allowbreak timestep}. The Professional rejects such inputs, so the
        branch is unreachable for Professional-aligned models, but it is a
        GNU-only semantic that can surprise a user. Either reject
        reset+timestep at the same index (matching the Professional) or document
        the branch prominently as a GNU extension and add a warning.

  \item[(f4) \code{ELEMENT\_\allowbreak MATRIX\_\allowbreak DELETE} dead record.]
        \code{ELEMENT\_\allowbreak MATRIX\_\allowbreak DELETE} and its allocation in \code{top.cc}/
        \code{database.cc} remain registered but are now unused (the
        Professional has no equivalent). Retire the record and its allocation,
        and the enum slot if a renumbering of the database enum is acceptable
        (the headers \code{tochnog.h}/\code{tochnog-mod.h} must stay in sync).

  \item[(f5) \code{refine\_\allowbreak l.cc} unreachable BAR3 branch.] A second copy-paste
        defect in \code{refine\_\allowbreak locally()}: the side-node list inside the
        \code{for (iside=...)} loop has \code{else if (name==-BAR2)} where
        \code{-BAR3} was intended (\code{refine\_\allowbreak l.cc:249}), so the branch is
        unreachable and the \code{-BAR3} local-refinement path would use a stale
        \code{ilist}. It is not exercised by any corpus test. Change the branch
        to \code{-BAR3} and add a focused \code{-bar3} local-refinement test.

  \item[(f6) Band width derived from \code{NODE\_\allowbreak NODE}.] The Bi-CG direct-LU
        retry band width is derived from the \code{NODE\_\allowbreak NODE} connectivity with
        an identity ordering (\code{(1+diff)*nprinc}). It is correct but
        conservative; a tighter exact bound could be derived from the element
        matrix unknown pairs, or a debug-build bounds assertion could be added in
        the band fill to catch any future out-of-band write.

  \item[(f7) \code{control\_\allowbreak eigen} is GNU-only.] \code{control\_\allowbreak eigen} /
        \code{node\_\allowbreak eigen} have no Professional equivalent (neither the
        25-10-2023 nor the 02-08-2026 binary contains the strings), so the eigen
        path has no convergence reference; the \code{ddc9a5b} fix restored the
        original GNU semantics, not a Professional alignment. Keep it documented
        as GNU-only and add an explicit note in the manual-user page.
\end{description}

\section{Recommendation}

The performance cost \textbf{(a)} is already addressed. Highest value for the
least risk among the remaining proposals: \textbf{(c)} (a clean build gate) and
\textbf{(d)} (a reproducible corpus). The principled long-term fix is
\textbf{(b)/(f2)}, but it depends on the shared-node tie consistency and is a
corpus-wide refactor.

% ===========================================================================
\appendix

\chapter{Corpus reconstruction method}
\label{app:reconstruction}

\code{tochnog} dumps the whole external database at the end of a run
(\code{print.cc} / \code{print\_\allowbreak db.cc}, \code{data\_\allowbreak number=-DATABASE}). The dump
mixes the parsed input records (canonicalised, comment-free, reordered) with
records the run itself produced, plus a few lost scalars.

\begin{enumerate}
  \item \textbf{Parser grounding.} The input grammar was read from
        \code{input.cc} (the init-section keyword list and the data-part
        \code{db\_\allowbreak number()} lookup). Input records are stored verbatim; the dump
        line format (\code{keyword [index] values}) already matches what the
        parser expects, so kept lines are emitted as-is.
  \item \textbf{Derived drop-list.} 103 record names are dropped. The core is
        the 45 names that appear in the 41 known \code{.dat}/\code{.dbs} pairs in
        \emph{no} \code{.dat} (state/results such as \code{node\_\allowbreak dof*},
        \code{element\_\allowbreak dof*}, \code{node\_\allowbreak rhside*}, \code{node\_\allowbreak start\_\allowbreak refined},
        \code{element\_\allowbreak matrix*}, \code{dof\_\allowbreak *}, \code{time\_\allowbreak *},
        \code{global\_\allowbreak solver\_\allowbreak *}, \code{post\_\allowbreak *}, \code{tochnog\_\allowbreak version},
        \code{Size}, \code{Total}), extended with the run-produced records seen
        in the wider dump set (\code{element\_\allowbreak beam\_\allowbreak *}, \code{element\_\allowbreak interface\_\allowbreak *},
        \code{element\_\allowbreak mass/\allowbreak volume}, \code{element\_\allowbreak geometry},
        \code{crack\_\allowbreak stressintensityfactor}, \code{inverse\_\allowbreak *}). Output provenance
        was confirmed in the source.
  \item \textbf{Initial-state recovery.} \code{node\_\allowbreak dof\_\allowbreak start\_\allowbreak refined} is the
        dump's copy of the start-of-run DOF (it holds the input \code{node\_\allowbreak dof}).
        It is re-emitted as \code{node\_\allowbreak dof}, otherwise hypoplasticity tests abort
        with ``Illegal porosity''. For full-mesh-generation tests
        (\code{control\_\allowbreak mesh\_\allowbreak macro}, \code{control\_\allowbreak mesh\_\allowbreak generate\_\allowbreak beam},
        \code{control\_\allowbreak mesh\_\allowbreak generate\_\allowbreak truss}) the generated
        \code{node}/\code{element}/\code{element\_\allowbreak group}/\code{node\_\allowbreak start\_\allowbreak refined}
        are dropped instead.
  \item \textbf{Lost scalars---derived, not invented.} The init-section
        \code{number\_\allowbreak of\_\allowbreak integration\_\allowbreak points} is stored bare in the dump (its
        value is never written). The value is recovered from the archived
        \code{element\_\allowbreak dof} record length and re-inserted into the reconstructed
        input.
\end{enumerate}

\section*{Derivation of the lost integration-point count}

In \code{elem.cc} the element state vector is allocated with a fixed stride

\begin{center}
\code{mnolnuknwn = npointmax * nuknwn}
\end{center}

and \code{db(ELEMENT\_\allowbreak DOF, ..., mnolnuknwn, ...)} dumps exactly that many values.
Here \code{npointmax} is the input \code{number\_\allowbreak of\_\allowbreak integration\_\allowbreak points} and
\code{nuknwn} is the number of element unknowns per integration point (a
function of the \code{derivatives} init set and \code{ndim} only, independent of
\code{npointmax}). The dump also contains the \code{dof\_\allowbreak label} record, whose
length is \code{nuknwn - 1} (the label array has one unused trailing slot).
Therefore:

\begin{center}
\code{nuknwn = len(dof\_\allowbreak label) + 1}\\
\code{number\_\allowbreak of\_\allowbreak integration\_\allowbreak points = len(element\_\allowbreak dof) /\allowbreak  nuknwn}
\end{center}

Both inputs are read from the archived \code{.dbs}; no fresh run is required and
no value is invented. The result is accepted only when the division is exact and
the element types are consistent with the resulting count.

\textbf{Cross-check with a fresh probe.} The reconstructed \code{.dat} (which
lacks the record) is re-run with an explicit
\code{number\_\allowbreak of\_\allowbreak integration\_\allowbreak points K} ($K=27$). Then
\code{len(element\_\allowbreak dof) = K * nuknwn}, so \code{nuknwn} can be measured on the
current binary. The two estimates agree on \textbf{72/72} tests, and the derived
counts agree on \textbf{72/72}.

\textbf{Validation on the known pairs.} 11 of the 41 known \code{.dat}/\code{.dbs}
pairs carry the record; for these the true value is known, and the archived-only
derivation reproduces it exactly in \textbf{11/11} cases (e.g.\ \code{gforce7}:
true 9, \code{dof\_\allowbreak label} 71, \code{nuknwn} 72, \code{element\_\allowbreak dof} 648, derived 9;
\code{msf\_\allowbreak cant3d\_\allowbreak hex27}: true 27, \code{dof\_\allowbreak label} 104, \code{nuknwn} 105,
\code{element\_\allowbreak dof} 2835, derived 27). The remaining 30 known pairs have no
\code{number\_\allowbreak of\_\allowbreak integration\_\allowbreak points} record in either file; applying the formula
there returns the archived binary's \emph{default} count, not an input value, so
those are correctly left untouched.

\section*{Caveats}

\begin{itemize}
  \item The archived \code{.dbs} were produced on 2026-08-29; the current binary
        contains many later physics commits, so even a perfect reconstruction
        need not reproduce the archived numbers.
  \item Records the run overwrites with its final value cannot be recovered as
        input (\code{node\_\allowbreak support\_\allowbreak edge\_\allowbreak normal\_\allowbreak force}). \code{node\_\allowbreak dof} is handled
        through \code{node\_\allowbreak dof\_\allowbreak start\_\allowbreak refined}.
  \item Inputs that are never echoed (\code{control\_\allowbreak solver}, \code{include},
        \code{support\_\allowbreak edge\_\allowbreak normal\_\allowbreak element\_\allowbreak node}, Abaqus/Gmsh/Feflow import
        records) are lost.
  \item \code{dbmeth21}/\code{dbmeth22} are \code{-size\_\allowbreak tot} dumps and contain no
        input.
  \item The derivation requires the \code{dof\_\allowbreak label} record; all 72 targets and
        all 41 known pairs have it. If a future dump lacks it, fall back to the
        K-probe cross-check.
\end{itemize}

\chapter{Reproducibility notes}
\label{app:repro}

\textbf{Corpus availability.} The sfnet 197
(\code{external-downloads/\allowbreak sfnet/\allowbreak extracted/\allowbreak test}) and validation-suite 233
(\code{validation-suite/\allowbreak test-2014}) directories are untracked/gitignored. The
honest denominators are 197 and 233; ``363'' was ephemeral and ``\textasciitilde207''
was a PASS count.

\textbf{Runner.} The corpus runner is a \code{/\allowbreak tmp} shell script
(\code{corpus-run.sh} / \code{run-corpus.sh} / \code{final-corpus-run.sh}
variants) with ad-hoc classification output. Target recalibration is manual: the
current value is measured, the \code{target\_\allowbreak value} line is edited by hand, and
the change is recorded verbatim in the docs (because git does not track the
corpus).

\textbf{Build and run discipline.} All measurements used clean
\code{build\_\allowbreak safe.sh} builds with \code{LC\_\allowbreak ALL=C} and a 45\,s per-test timeout.
Only clean builds are reliable on the current filesystem; incremental builds mix
old/new enums and corrupt the binary. The headers \code{tochnog.h} and
\code{tochnog-mod.h} must stay in sync (same enum, same order).

\textbf{Internal suite.} The \code{rc}-based internal suite is 16/16 and is the
primary gate. The \code{build\_\allowbreak safe.sh} full-mode file-check layer currently exits
non-zero on stale expected outputs (see improvement~(c)); its FALLO list is
documented to be identical to the baseline, i.e.\ it is not a regression.

\textbf{Reference binaries.} The behavioural reference is the Tochnog
Professional binary (build of 02-08-2026, with the 25-10-2023 build for older
paths), together with a build of the reference GNU 2014 sources.

\textbf{Documentation convention.} Each change updated
\code{SEGUIMIENTO-\allowbreak CONVERGENCIA.md} and the relevant
\code{manual-user}/\code{manual-developer} pages in the same work unit, per the
project's mandatory convention.

\end{document}
